% Guia do professor — Capítulo 6: Nuvens
\section{Capítulo 6 --- Nuvens}

\subsection*{Visão geral e ritmo}

O capítulo mais visual do curso: a taxonomia das nuvens (linguagem
operacional) e os três mecanismos de saturação (física), com nevoeiros como
estudo de caso e a geometria fractal como sobremesa para físicos.
\textbf{Ritmo sugerido: 2 aulas de 2\,h} — pode comprimir para 1 aula e
meia se o semestre apertar (a classificação anda rápido com fotos).

\begin{itemize}
  \item \textbf{Aula 1} — gêneros com fotos (15\,min bem gastos),
        gramática dos nomes, os três caminhos para saturar (com a dedução
        da convexidade), bafo e contrails. Caso~1 (10\,min).
  \item \textbf{Aula 2} — nevoeiros (tabela + EDO do nevoeiro de
        radiação), organização (ruas, ondas, células), fractais, oktas.
        Casos~2--4 (20\,min; o Caso~3 rende).
\end{itemize}

\subsection*{Objetivos de aprendizagem}

\begin{enumerate}
  \item nomear os 10 gêneros e inferir o mecanismo (estável × convectivo)
        do aspecto;
  \item explicar os três caminhos para a saturação na curva $e_s(T)$ e
        provar por que mistura satura (convexidade);
  \item classificar nevoeiros pelo mecanismo e estimar tempo de formação do
        nevoeiro de radiação;
  \item relacionar padrões de organização com $z_i$ e com $N$ (Cap.~5);
  \item explicar o que significa dimensão fractal $D \simeq 1{,}35$ e como
        se mede;
  \item usar oktas e teto na linguagem METAR.
\end{enumerate}

\subsection*{Pré-requisitos a reativar}

Curva $e_s(T)$ e LCL (Cap.~4); $N$ de Brunt--Väisälä e camada de mistura
(Cap.~5); balanço radiativo noturno (Cap.~2). Fractais: nenhum
pré-requisito — introduzir do zero é parte da graça.

\subsection*{Armadilhas conceituais frequentes}

\begin{itemize}
  \item \textbf{``Nuvem é vapor''}: nuvem é água \emph{líquida/gelo};
        vapor é invisível. Corrigir logo com o bafo.
  \item \textbf{``Misturar ar só faz média''}: a média é na corda, mas a
        saturação é na curva — a convexidade surpreende até veteranos.
  \item \textbf{Nevoeiro como ``nuvem baixa'' sem física}: cada tipo tem um
        mecanismo distinto; a tabela de 5 tipos organiza.
  \item \textbf{Contrail como ``fumaça''}: é nuvem de mistura; a fuligem só
        fornece núcleos (gancho para o Cap.~7).
  \item \textbf{Perímetro fractal}: alunos esperam convergência ao refinar
        a régua; a divergência $\ell^{1-D}$ é contraintuitiva — o exercício
        T4 fixa.
\end{itemize}

\subsection*{Demonstração de baixo custo}

Nuvem na garrafa: garrafa PET com um dedo d'água + fumaça de palito de
fósforo; apertar e soltar — expansão adiabática forma a nuvem no ato
(mecanismo 2 + núcleos do Cap.~7). Complemento: vídeo timelapse de nevoeiro
de radiação se formando em vale.

\subsection*{Problemas sugeridos do livro (exercícios do Cap.~6)}

Aquecimento: identificação de gêneros a partir de descrições; altura de
base de cumulus via LCL. Aprofundamento: problemas de mistura
(bafo/contrail); tempo de formação de nevoeiro. Síntese: ``por que noites
nubladas não dão nevoeiro de radiação?''.
\emph{Confira a numeração na sua cópia (v1.02b).}

\subsection*{Uso do notebook \texttt{cap06\_nuvens.ipynb}}

\begin{center}
\begin{tabular}{lll}
\hline
Caso & Conteúdo & Momento sugerido \\
\hline
1 & Diagrama de mistura: bafo e contrail em $e_s(T)$ & Aula 1, após mistura \\
2 & EDO do nevoeiro de radiação (forma e dissipa) & Aula 2, após nevoeiros \\
3 & Nuvens fractais sintetizadas; $D$ por contagem de caixas & Aula 2, fractais \\
4 & $\lambda$ de ondas e ruas a partir da sondagem & Aula 2, organização \\
\hline
\end{tabular}
\end{center}

O Caso~3 é o mais ``físico computacional'' do curso até aqui (síntese
espectral + contagem de caixas) — bom para apresentar como mini-projeto.

\subsection*{Exercícios complementares e soluções}

Nas notas: T1--T5 (convexidade e supersaturação de misturas, critério de
Appleman, dedução do $t_{fog}$, aritmética fractal, espaçamentos de
organização) e N1--N4. Soluções em \texttt{cap06\_solucoes.ipynb}
(\emph{não distribuir}): T1 com \texttt{sympy} (segunda derivada de
Clausius--Clapeyron), N3 com o estudo $D(\beta)$ completo. Pares:
T1$\leftrightarrow$N1, T3$\leftrightarrow$N2, T4$\leftrightarrow$N3.

\subsection*{Conexões com o resto do curso}

LCL/base de cumulus $\to$ Cap.~4; estável × convectivo e $N$ $\to$ Cap.~5;
núcleos de condensação e crescimento de gotas $\to$ Cap.~7; oktas/teto
$\to$ METARs (Cap.~9); nevoeiro de radiação $\to$ camada limite noturna
(Cap.~18); células abertas/fechadas $\to$ interação ar-mar (Cap.~21).
